\documentclass[12pt]{article}
\usepackage[utf8]{inputenc}
\usepackage[T1]{fontenc}
\usepackage{hyperref}
\title{North Fali Language Resources for K-12 Curriculum Development}
\author{Compiled by Claude Sonnet 5.5 (claude-sonnet-5-5) for Andr\'as Kornai}
\date{}
\begin{document}
\maketitle

\section{Introduction}
North Fali (ISO 639-3 fll) is a small Atlantic-Congo language of northern Cameroon, with roughly 16,000 speakers according to Wikipedia. Its varieties are listed there as Dourbeye, Bossum and Bv\textschwa ri; it should not be confused with South Fali (including Fali-Tinguelin). Wikipedia notes that speakers were shifting to Adamawa Fulfulde by 1982. Very little material exists, and none found is designed for schools.

\section{Primary Educational Resources}
No primers, textbooks or school materials in North Fali were found in web searches.

\section{Linguistic and Reference Materials}
\begin{itemize}
\item Raija Kramer, \emph{Die Sprache der Fali in Nordkamerun} (2014), a German-language grammar published by K\"oppe: \url{https://www.koeppe.de/titel_die-sprache-der-fali-in-nordkamerun}
\item SIL archive record of Ennulat, \emph{Outline of the grammar of Fali} (1971), tagged to North Fali; not found to be available online: \url{https://www.silcam.org/resources/archives/64940}
\item Overview of the Fali languages of Cameroon, with further reading (including Sweetman 1981, a comparative study of Fali dialects): \url{https://en.wikipedia.org/wiki/Fali_languages_(Cameroon)}
\end{itemize}

\section{Religious and Community Resources}
Joshua Project has a people-group page for North Fali in Cameroon: \url{https://joshuaproject.net/people_groups/13981}. No scripture or community organisation was found.

\section{Digital and Audio Resources}
No audio, video, corpus or software resources were found.

\section{Curriculum Development Context}
The language is not known to be taught in schools, and no orthography or pedagogical materials were found. A curriculum would need to begin from the Kramer grammar and new fieldwork with the community; Fulfulde and French dominate regional schooling.
\end{document}